\chapter*{Guide institutionnel int\'egr\'e du stage PFE}
\addcontentsline{toc}{chapter}{Guide institutionnel int\'egr\'e du stage PFE}

\begin{guidebox}
Cette partie transforme le template en support de lecture du r\'ef\'erentiel PFE. Elle est visible uniquement lorsque le mode guide est activ\'e. Pour la version finale destin\'ee au d\'ep\^ot, d\'esactivez-la en rempla\c{c}ant \texttt{\textbackslash showguidetrue} par \texttt{\textbackslash showguidefalse} dans \texttt{config.tex}.
\end{guidebox}

\section*{1. Objectifs du stage de PFE}
\addcontentsline{toc}{section}{1. Objectifs du stage de PFE}

Le stage de Projet de Fin d'\'Etudes constitue une situation de synth\`ese et de professionnalisation. Il a pour finalit\'e de permettre \`a l'\'el\`eve-ing\'enieur de d\'emontrer qu'il est capable d'aborder un probl\`eme complexe d'ing\'enierie, d'en analyser les enjeux, de proposer une d\'emarche rigoureuse, puis de d\'evelopper et de valider une solution argument\'ee dans un contexte r\'eel ou proche du r\'eel. Le rapport n'est donc pas un simple compte rendu d'activit\'e, mais une preuve acad\'emique de comp\'etences. \`A minima, le stage doit conduire l'\'etudiant \`a :

\begin{itemize}[leftmargin=1.2cm]
\item analyser le contexte, la probl\'ematique, les besoins, les contraintes et le p\'erim\`etre du sujet ;
\item mobiliser un \`etat de l'art, des normes, des outils, des m\'ethodes ou des solutions existantes de mani\`ere critique ;
\item justifier des choix techniques, scientifiques, m\'ethodologiques ou \`economiques ;
\item concevoir, impl\'ementer, mod\'eliser, exp\'erimenter ou int\'egrer une solution ;
\item valider les r\'esultats \`a l'aide d'un protocole explicite et d'indicateurs pertinents ;
\item faire preuve d'autonomie, de professionnalisme, d'int\'egrit\'e scientifique et de capacit\'e \`a communiquer clairement.
\end{itemize}

\guided{Dans l'introduction g\'en\'erale, reformulez le sujet sous forme de probl\'ematique et annoncez des objectifs mesurables. Dans la conclusion, revenez explicitement sur le niveau d'atteinte de ces objectifs.}

\section*{2. Ce que le rapport doit d\'emontrer}
\addcontentsline{toc}{section}{2. Ce que le rapport doit d\'emontrer}

Le lecteur doit pouvoir reconstituer la cha\^ine logique suivante : \textbf{probl\`eme ou besoin $\rightarrow$ exigences $\rightarrow$ m\'ethode $\rightarrow$ solution propos\'ee $\rightarrow$ validation $\rightarrow$ r\'esultats $\rightarrow$ conclusion}. En pratique, cela signifie que le rapport doit pr\'esenter de mani\`ere tra\c{c}able :

\begin{itemize}[leftmargin=1.2cm]
\item le contexte et le positionnement du projet ;
\item les besoins fonctionnels et non fonctionnels, les contraintes et les hypoth\`eses ;
\item les choix de conception et leurs justifications ;
\item la r\'ealisation ou l'exp\'erimentation effective ;
\item un plan de tests ou un protocole de validation ;
\item une discussion critique : limites, \`ecarts, menaces sur la validit\'e, perspectives.
\end{itemize}

\guided{Un bon rapport de PFE n'est pas seulement descriptif. Il doit montrer comment et pourquoi certaines d\'ecisions ont \`et\'e prises, puis comment leur pertinence a \`et\'e v\'erifi\'ee.}

\section*{3. Modalit\'es d'\'evaluation recommand\'ees}
\addcontentsline{toc}{section}{3. Modalit\'es d'\'evaluation recommand\'ees}

L'\'evaluation du PFE doit rester lisible, homog\`ene et align\'ee sur les comp\'etences vis\'ees. Le tableau suivant peut \^etre conserv\'e dans le template pour que l'\'etudiant comprenne, d\`es la phase de r\'edaction, ce qui sera observ\'e par le jury.

\begin{longtable}{>{\RaggedRight\arraybackslash}p{4.3cm} >{\RaggedRight\arraybackslash}p{8cm} >{\centering\arraybackslash}p{2cm}}
\toprule
\textbf{Crit\`ere} & \textbf{Indicateurs observables} & \textbf{Poids indicatif} \\
\midrule
\endfirsthead
\toprule
\textbf{Crit\`ere} & \textbf{Indicateurs observables} & \textbf{Poids indicatif} \\
\midrule
\endhead
Qualit\'e du cadrage et de l'analyse & Probl\'ematique, objectifs, besoins, p\'erim\`etre, positionnement initial. & 10\% \\
\addlinespace
\'Etat de l'art et usage des sources & Pertinence, diversit\'e, esprit critique, qualit\'e des citations et du positionnement. & 10\% \\
\addlinespace
M\'ethodologie et rigueur & Choix justifi\'es, d\'emarche structur\'ee, tra\c{c}abilit\'e, logique d'ing\'enierie. & 15\% \\
\addlinespace
Conception et r\'ealisation & Architecture, impl\'ementation, int\'egration, qualit\'e technique de la solution. & 20\% \\
\addlinespace
Validation et discussion & Protocoles, r\'esultats, analyse critique, comparaison aux objectifs. & 15\% \\
\addlinespace
Qualit\'e du rapport & Structure, clart\'e, style, figures, tableaux, coh\'erence formelle. & 10\% \\
\addlinespace
Soutenance orale & Clart\'e, ma\^itrise du sujet, gestion du temps, qualit\'e des supports. & 10\% \\
\addlinespace
R\'eponses aux questions et recul critique & Argumentation, autonomie, lucidit\'e sur les limites et les perspectives. & 10\% \\
\bottomrule
\end{longtable}

\guided{Utilisez cette grille comme un outil d'auto-contr\^ole. Avant le d\'ep\^ot, v\'erifiez que chaque crit\`ere est visible dans le rapport ou pr\'epar\'e pour la soutenance.}

\section*{4. R\'ef\'erences bibliographiques : exigences et bonnes pratiques}
\addcontentsline{toc}{section}{4. R\'ef\'erences bibliographiques : exigences et bonnes pratiques}

La bibliographie fait partie des indicateurs majeurs de la qualit\'e scientifique d'un rapport. Toute id\'ee, donn\'ee, figure, tableau, extrait de code, formule ou m\'ethode emprunt\'e(e) doit \^etre attribu\'e(e) correctement. Le rapport doit respecter les principes suivants :

\begin{itemize}[leftmargin=1.2cm]
\item utiliser un style homog\`ene ; le style \textbf{IEEE} est recommand\'e pour les fili\`eres d'ing\'enierie ;
\item s'assurer que toute source cit\'ee dans le texte figure dans la bibliographie et r\'eciproquement ;
\item privil\'egier les articles index\'es, actes de conf\'erence reconnus, ouvrages acad\'emiques, normes, brevets et rapports techniques fiables ;
\item viser au moins \textbf{\MinReferences{} r\'ef\'erences} pour un niveau master / cycle ing\'enieur, dont \IndexedRatio{} de sources scientifiques index\'ees lorsque cela est applicable ;
\item \`eviter Wikipedia, les blogs non sp\'ecialis\'es, les contenus commerciaux non v\'erifiables et les sites sans auteur identifiable ;
\item ajouter la date de consultation pour les ressources web.
\end{itemize}

\subsection*{4.1 Comment citer dans le texte}

En style IEEE, les citations apparaissent entre crochets num\'erot\'es : par exemple \texttt{[1]}, \texttt{[2]} ou \texttt{[3], [4]}. La citation doit appara\^itre au plus pr\`es de l'id\'ee concern\'ee. Une paraphrase doit \^etre cit\'ee tout autant qu'une citation directe.

\begin{quote}
Exemple : \textit{Les travaux r\'ecents montrent que la robustesse du syst\`eme d\'epend fortement de la qualit\'e des donn\'ees d'apprentissage et du protocole de validation} \cite{example_journal_2024}. \\
Exemple : \textit{L'architecture retenue est coh\'erente avec les pratiques d\'ecrites dans la litt\'erature technique} \cite{example_conference_2023,example_standard}.
\end{quote}

\subsection*{4.2 Exemples de r\'ef\'erences au format IEEE}

\begin{itemize}[leftmargin=1.2cm]
\item \textbf{Article de revue :} [1] A. Auteur et B. Auteur, \enquote{Titre de l'article,} \textit{Nom de la revue}, vol. X, no. Y, pp. xx--yy, ann\'ee.
\item \textbf{Conf\'erence :} [2] A. Auteur, \enquote{Titre du papier,} in \textit{Proc. Nom de la conf\'erence}, ville, pays, ann\'ee, pp. xx--yy.
\item \textbf{Ouvrage :} [3] A. Auteur, \textit{Titre du livre}, Xe \'ed. ville : \'editeur, ann\'ee.
\item \textbf{Norme :} [4] \textit{Titre de la norme}, organisme, r\'ef\'erence, ann\'ee.
\item \textbf{Site web :} [5] A. Auteur ou organisme, \enquote{Titre de la ressource,} site, ann\'ee. [En ligne]. Disponible : URL. [Consult\'e le : jj mois aaaa].
\end{itemize}

\subsection*{4.3 Comment remplir \texttt{references.bib}}

Le fichier \texttt{references.bib} fourni avec ce template permet \`a l'\'etudiant de g\'erer ses sources proprement avec BibTeX. La d\'emarche minimale est la suivante :

\begin{enumerate}[leftmargin=1.2cm]
\item ajouter une entr\'ee bibliographique dans \texttt{references.bib} ;
\item citer la source dans le texte avec \texttt{\textbackslash cite\{cle\}} ;
\item compiler le projet selon la s\'equence standard d'Overleaf ;
\item v\'erifier que toutes les r\'ef\'erences sont correctes, compl\`etes et homog\`enes.
\end{enumerate}

\subsection*{4.4 Exemples BibTeX \`a recopier et adapter}

\begin{lstlisting}[numbers=none]
@article{nomArticle2025,
  author  = {A. Nom and B. Nom},
  title   = {Titre de l'article},
  journal = {Nom de la revue},
  volume  = {12},
  number  = {3},
  pages   = {101--115},
  year    = {2025}
}

@inproceedings{nomConf2024,
  author    = {A. Nom and B. Nom},
  title     = {Titre du papier},
  booktitle = {Proceedings of ...},
  pages     = {55--60},
  year      = {2024}
}

@book{nomLivre2023,
  author    = {A. Nom},
  title     = {Titre du livre},
  publisher = {Editeur},
  year      = {2023}
}

@misc{nomWeb2026,
  author = {{Nom de l'organisme}},
  title  = {Titre de la ressource web},
  year   = {2026},
  url    = {https://exemple.org},
  note   = {[En ligne]. Consulte le : 15 avril 2026}
}
\end{lstlisting}

\guided{Une bibliographie longue mais faible n'am\'eliore pas la qualit\'e du rapport. Le jury valorise davantage une s\'election pertinente, r\'ecente, scientifiquement cr\'edible et effectivement mobilis\'ee dans l'analyse.}

\section*{5. Conseils d'utilisation p\'edagogique du template}
\addcontentsline{toc}{section}{5. Conseils d'utilisation p\'edagogique du template}

Il est recommand\'e que l'\'etudiant lise cette partie avant de commencer la r\'edaction, puis y revienne \`a trois moments-cl\'es :

\begin{itemize}[leftmargin=1.2cm]
\item au d\'emarrage, pour comprendre la finalit\'e acad\'emique du stage ;
\item \`a mi-parcours, pour v\'erifier l'alignement entre objectifs, chapitres et preuves produites ;
\item avant le d\'ep\^ot, pour contr\^oler la conformit\'e du rapport avec les attentes institutionnelles.
\end{itemize}

\begin{guidebox}
Suggestion pratique pour l'encadrant : conserver cette partie dans la version de travail partag\'ee avec l'\'etudiant, puis la masquer dans la version finale de d\'ep\^ot. Le template joue ainsi un double r\^ole : \textbf{mod\`ele de rapport} et \textbf{support de lecture du r\'ef\'erentiel PFE}.
\end{guidebox}
